\documentclass[12pt]{article}
\usepackage{fullpage,hyperref}\setlength{\parskip}{3mm}\setlength{\parindent}{0mm}
\begin{document}

\begin{center}\bf
Homework 9. Due by 11:59pm on Sunday 11/15.

Linux and the command line.

\end{center}

Linux is the dominant environment for scientific computing. For example, all maior supercomputers run some variant of Linux (\url{https://en.wikipedia.org/wiki/TOP500}). As another example, most cloud servers are built on Linux, and Linux is therefore dominant for data science applications that involve cloud computing. At University of Michigan, the main resource for high performance computing is the Great Lakes Linux cluster. It should be apparent that Linux skills are useful for a research statistician, as soon as your data analysis or simulation study is too large for a laptop.

Linux expertise in this class ranges from novice to expert. Our goal is to advance our understanding and share knowledge. Write brief answers to the following questions, by editing the tex file available at \url{https://github.com/ionides/810f26}, and submit the resulting pdf file via Canvas.

\begin{enumerate}
  
\item To what extent do you agree or disagree with the opinions at
\url{https://valohai.com/blog/command-line-for-data-science/}?
Specifically,

(i) ``most data scientists are on UNIX-based systems these days.''

(ii) Five reasons why command line use is critical for data science: speed, agnosticism, automation, extensibility, lack of other options.


YOUR ANSWER HERE.

\item  Unless you vigorously disagree with the assertions in Question 1, if you have not yet had much experience working with the command line, you should be motivated to change that! The instructor of this course endorses the value of working with the command line, so even if you are deeply skeptical about the value of learning the ancient art of command line computing, please put that to one side for the rest of the course.

If you are new to Linux, or if your Linux skills are limited to a handful of commands, read the introduction to command line Linux at

{\small\hspace{-1cm}\url{https://ubuntu.com/desktop/docs/en/latest/tutorial/the-linux-command-line-for-beginners}}

If you have a Mac, try out the commands on a Terminal app, which runs a version of Unix that works identically to Linux for many everyday purposes. If you run Windows, try out some commands on the Windows subsystem for Linux

\url{https://docs.microsoft.com/en-us/windows/wsl/}

You can also log in to the UM Linux login service, using SSH Secure Shell. For example from a Mac terminal, type

\texttt{ssh your\_uniqname@login.itd.umich.edu}

Optionally, if you find it an effective way to practice, play the terminus adventure at

\url{https://web.mit.edu/mprat/Public/web/Terminus/Web/main.html}

If you are already comfortable with the command line, look through the materials on the MIT ``Missing Semester'' course at \url{https://missing.csail.mit.edu/} until you find something new and interesting.

Report briefly on which activity you chose and what you have learned.

YOUR ANSWER HERE.


\item In what ways has the arrival of agentic AI increased and decreased the role of the command-line for scientific computing (broadly interpreted to include statistics)?

YOUR ANSWER HERE.

\item If you are a relatively experienced Linux user, share some words of advice for beginners. How and why did you get started with Linux?  

YOUR ANSWER HERE.

\end{enumerate}
\end{document}
